\chapter{Capability Models and Policy Translation: SCM}
\label{ch:scm}
\draftnote{Target length: about 6 pages. Source: \cite{basile2026scm} (\PB).}

\section{The Configuration Problem}
\label{sec:scm-problem}
\draftnote{Heterogeneous security controls, vendor lock-in, and configuration as a
historically manual and error-prone activity.}

\section{The Information Model and the Data Model}
\label{sec:scm-models}
\draftnote{Conditions, actions and events; policies with condition evaluation, resolution
strategies and default actions.}

\section{The Three-Rung Ladder}
\label{sec:scm-ladder}
\draftnote{High-Level Policy $\to$ Medium-Level Policy $\to$ Low-Level Configuration, and
what each rung is expressive enough to say. A TikZ figure of the ladder fits here.}

\section{Model-Driven Translation}
\label{sec:scm-translation}
\draftnote{Where the descent rules actually live: the translation and policy-detail
association classes, and the decorator pattern that keeps device specifics out of the
translator code.}

\section{How the Translation Is Validated}
\label{sec:scm-validation}
\draftnote{A refinement engine implemented as forward reasoning, expert manual validation
of generated artifacts, and attack-simulation scripts. State plainly: no proof obligation is
formulated for any rung of the ladder.}

\section{Assessment}
\label{sec:scm-assessment}
\draftnote{The most explicit abstraction ladder of the four approaches, and the weakest
guarantee that each step down it is sound. Fill the SCM rows of
\cref{tab:models,tab:evidence}.}
